% !TeX program = tectonic
\documentclass[11pt,a4paper]{article}
\usepackage{fontspec}
\usepackage[russian]{babel}
\babelfont{rm}[Language=Default]{Liberation Serif}
\babelfont[Language=Default]{sf}{Carlito}
\babelfont[Language=Default]{tt}{DejaVu Sans Mono}
\usepackage{amsmath,amssymb,amsthm}
\usepackage{geometry}
\geometry{a4paper, top=2.4cm, bottom=2.4cm, left=2.4cm, right=2.4cm}
\usepackage{booktabs,tabularx,enumitem,xcolor,tcolorbox}
\tcbuselibrary{breakable,skins}
\definecolor{linkblue}{HTML}{1F4E79}
\definecolor{statgreen}{HTML}{2F6B3A}
\definecolor{goldacc}{HTML}{8B7E5A}
\usepackage[colorlinks=true,linkcolor=linkblue,urlcolor=linkblue,unicode=true]{hyperref}
\theoremstyle{plain}
\newtheorem{theorem}{Теорема}
\newtheorem{lemma}[theorem]{Лемма}
\newtheorem{proposition}[theorem]{Предложение}
\newtheorem{corollary}[theorem]{Следствие}
\theoremstyle{definition}
\newtheorem{definition}[theorem]{Определение}
\theoremstyle{remark}
\newtheorem{remark}[theorem]{Замечание}
\newtcolorbox{statusbox}[1][]{colback=statgreen!5,colframe=statgreen!75,
  title={\textbf{Сводка доказанного:} #1},fonttitle=\small,boxrule=0.7pt,arc=2pt,breakable}
\newtcolorbox{databox}[1][]{colback=goldacc!6,colframe=goldacc!80,
  title={\textbf{Данные протокола:} #1},fonttitle=\small,boxrule=0.7pt,arc=2pt,breakable}
\newtcolorbox{verifybox}[1][]{colback=linkblue!5,colframe=linkblue!80,
  title={\textbf{Верификация:} #1},fonttitle=\small,boxrule=0.7pt,arc=2pt,breakable}
\widowpenalty=10000
\clubpenalty=10000
\setlength{\parskip}{0.35em}
\newcommand{\CC}{\mathbb{C}}
\newcommand{\RR}{\mathbb{R}}
\newcommand{\ZZ}{\mathbb{Z}}
\newcommand{\QQ}{\mathbb{Q}}
\newcommand{\Ga}{\Gamma}
\newcommand{\Bfun}{\mathrm{B}}
\newcommand{\im}{\operatorname{Im}}
\newcommand{\re}{\operatorname{Re}}
\newcommand{\lcm}{\operatorname{lcm}}
\newcommand{\SNF}{\operatorname{SNF}}
\newcommand{\Jac}{\operatorname{Jac}}
\newcommand{\rank}{\operatorname{rank}}
\newcommand{\Ind}{\operatorname{Index}}
\newcommand{\disc}{\operatorname{disc}}
\begin{document}
\begin{center}
{\LARGE\bfseries Тройка тора и вывод поправок}\\[0.4em]
{\large калибровочный стенд из первых принципов}\\[0.6em]
{\normalsize Исаев Исхак Хамзатович}\\[0.2em]
{\small Программа <<Динамический принцип>> --- лаборатория \texttt{hodge-laboratory}}\\[0.15em]
{\small Отдельная монография теоремы · Монография 7/16}
\end{center}
\vspace{0.4em}
{\color{linkblue}\hrule height 0.8pt}
\vspace{0.8em}

\section{Постановка}

Тор --- калибровочный стенд: его тройка $(4,1,4\pi^2)$ и все
производные величины вычисляются из первых принципов, без единой
внешней посылки. Именно на торе конвейер поправок виден целиком, и
любая ошибка конструкции проявилась бы при грубом счёте.

\begin{lemma}[тройка тора]\label{lem:torus}
Тор поставляет целочисленную тройку $(4,\,1,\,4\pi^2)$: порядок
поворота $n=4$, индекс носителя $B=1$, спектральный порог
$\lambda_0=4\pi^2$.
\end{lemma}

\begin{proof}
Поворот $x\mapsto ix$ имеет порядок $4$: $i^4=1$, $i^2=-1\ne1$.
Носитель --- сам тор: $H_2(T^2,\ZZ)\cong\ZZ$, $b_2=1$. Спектр
плоского лапласиана --- $|m\tau-n|^2$; в нормировке Tate порог
равен $4\pi^2$ --- универсальный первый член нормировки
(сертификат G1).
\end{proof}

\begin{theorem}[полный вывод поправок на торе]\label{th:toruscorr}
Из тройки $(4,1,4\pi^2)$ правила конвейера дают:
$\delta=\pi/4=0.7853981634$; голономия поворота $=i$ (отклонение
протокола $6.1\cdot10^{-17}$ --- машинная точность);
$\delta^2/2=0.3084251375$; $\gamma=\delta^4/1=0.3805042619$;
$\delta_{\mathrm{eff}}=\delta^5=0.2988473484$;
$\Delta_{Ch}^{\mathrm{torus}}=4\pi^2+\frac{\delta^2}2-\frac{\delta^5}1
=39.4879953935$.
\end{theorem}

\begin{proof}
Каждая величина --- по одному правилу: $\delta=\pi/N$
(спиновое двойное накрытие); $\delta^2/2$ --- главная фазовая
поправка; $\gamma=\delta^4/k$ при $k=1$; $\delta_{\mathrm{eff}}
=\delta\cdot\gamma$; объединение --- по формуле $\Delta_{Ch}$ при
$R=0$. Голономия: подъём поворота в спин-расслоение переставляет
листы, собственное значение на собственной форме равно $i$.
Все знаки воспроизводятся при $dps=35$ с нулевой невязкой
последнего разряда.
\end{proof}

\begin{remark}[эквивариантность шахматной сетки]
Для симметричной кодировки $\mu_4$-эквивариантность точная
(отклонение $0.000$ по сетке $4096\times4096$); для кодировки
$(7,22)$ --- нарушена ($1.000$). Контраст доказывает: поправки
геометричны --- они наследуют симметрию носителя.
\end{remark}

\begin{databox}[числа]
Семейство фаз $\delta_T=\pi/N$ воспроизведено для $N=2,\dots,8$;
таблицы $\gamma$ и $\delta_{\mathrm{eff}}$ для $k\in\{1,22\}$
построены полностью; кейс E6 исправлен: $\Delta=0.7990$
(не $0.5964$); на Клейне та же формула даёт $\Delta_{Ch}\approx3.4399$
при наблюдаемом $3.443$ --- отклонение $0.0905\%$.
\end{databox}

\begin{statusbox}[итог]
Тройка тора и все производные величины выведены из первых принципов;
конвейер на торе работает точно, что является калибровкой для всех
верхних ступеней.
\end{statusbox}

\begin{verifybox}[как воспроизвести]
\texttt{python3 laboratory.py} $\to$ <<Стенды $\to$ Тор>>; таблица
$\gamma,\delta_{\mathrm{eff}}$ --- пункт <<Таблицы>>;
диаграмма \texttt{fig\_torus.png}.
\end{verifybox}

\vfill
{\color{linkblue}\hrule height 0.6pt}
\vspace{0.5em}
{\small\color{gray}
Лицензия: индивидуальная исключительная лицензия Исаева Исхака Хамзатовича.
Все права на вычисления, формулы и тексты принадлежат автору.
Верификация: \texttt{python3 laboratory.py} (пункт <<Стенды>>/<<Сертификаты>>),
Lean: \texttt{verification/lean/HodgeLaboratory.lean}.
}
\end{document}
